\chapter{Całka Riemanna}
\label{ma2:ch:riemann-integral}

Podstawową konstrukcję całki Riemanna na przedziale,
równoważność sum Riemanna i Darboux, własności
algebraiczne oraz przykłady omówiono w analizie
matematycznej I:
\atlasref{AM1-29::ma1:ch:riemann-integral}{całka Riemanna}.
Tutaj przygotujemy jej rozszerzenie na prostokąty
w $\R^n$. Zagadnienia całek iterowanych, zamiany
zmiennych i całkowania po ogólniejszych zbiorach
rozwiniemy w rozdziale poświęconym całkom wielokrotnym.
Całka funkcji o wartościach w przestrzeni Banacha
będzie przedmiotem osobnego rozdziału.

\section{Podziały prostokąta i sumy Darboux}

Ustalmy $n\ge1$ i prostokąt domknięty
\[
 Q=\prod_{j=1}^n[a_j,b_j],\qquad a_j<b_j.
\]
Przez objętość prostokąta rozumiemy liczbę
$|Q|=\prod_{j=1}^n(b_j-a_j)$. Niech w każdym
przedziale współrzędnościowym będzie dany podział
$a_j=x_{j,0}<\cdots<x_{j,N_j}=b_j$.
Ich rodzina $P$ wyznacza komórki
\[
 Q_\alpha=\prod_{j=1}^n
 [x_{j,\alpha_j-1},x_{j,\alpha_j}],
 \qquad 1\le\alpha_j\le N_j.
\]
Komórki są domknięte, pokrywają $Q$ i mają parami
rozłączne wnętrza. Objętości spełniają
\[
 \sum_\alpha |Q_\alpha|
 =\prod_{j=1}^n
   \sum_{i=1}^{N_j}(x_{j,i}-x_{j,i-1})=|Q|.
\]
Równość wynika z rozdzielności mnożenia względem
skończonych sum i teleskopowania w każdej współrzędnej.
Przez średnicę podziału rozumiemy
$\delta(P)=\max_{j,i}(x_{j,i}-x_{j,i-1})$;
średnica euklidesowa każdej komórki nie przekracza
$\sqrt n\,\delta(P)$.

\begin{definicja}
\label{ma2:def:darboux-box}
Dla ograniczonej funkcji $f:Q\to\R$ definiujemy
\[
 \begin{aligned}
 m_\alpha&=\inf_{Q_\alpha}f,\qquad
 M_\alpha=\sup_{Q_\alpha}f,\\
 L(f,P)&=\sum_\alpha m_\alpha|Q_\alpha|,\qquad
 U(f,P)=\sum_\alpha M_\alpha|Q_\alpha|.
 \end{aligned}
\]
Całka dolna i górna Darboux to odpowiednio
$\underline I(f)=\sup_P L(f,P)$ oraz
$\overline I(f)=\inf_P U(f,P)$.
Funkcja jest całkowalna w sensie Riemanna na $Q$,
gdy $\underline I(f)=\overline I(f)$.
Wspólną wartość oznaczamy $I_Q(f)$ lub
$\int_Q f(x)\,\mathrm d\lambda_n(x)$ z zaznaczeniem,
że na tym etapie jest to całka Riemanna.
Zgodność z całką Lebesgue'a uzasadnimy dalej.
\end{definicja}

\begin{stwierdzenie}[Kryterium Darboux]
\label{ma2:prop:darboux-box}
Zachodzi $\underline I(f)\le\overline I(f)$.
Funkcja ograniczona $f$ jest całkowalna wtedy
i tylko wtedy, gdy dla każdego $\varepsilon>0$
istnieje podział $P$ taki, że
$U(f,P)-L(f,P)<\varepsilon$.
\end{stwierdzenie}
\begin{proof}
Rozdrobnienie podziału oznacza dodanie węzłów
w dowolnych współrzędnych. Każda nowa komórka
zawiera się w starej, więc jej infimum jest nie
mniejsze, a supremum nie większe niż odpowiednie
kresy w starej komórce. Suma objętości nowych
komórek zawartych w starej jest równa jej
objętości (na mocy wyżej wykazanej tożsamości).
Stąd suma dolna nie maleje, a górna nie rośnie.
Dla dwóch podziałów $P,T$ tworzymy wspólne
rozdrobnienie $R$ przez sumę ich węzłów w każdej
współrzędnej. Otrzymujemy
\[
 L(f,P)\le L(f,R)\le U(f,R)\le U(f,T).
\]
Przejście do kresów daje pierwszą tezę.
Wszystkie kresy są skończone, bo $f$ jest ograniczona,
a suma objętości komórek wynosi $|Q|$.
Jeżeli całki dolna i górna są równe $I$, wybieramy
$P,T$ z $L(f,P)>I-\varepsilon/2$ i
$U(f,T)<I+\varepsilon/2$. Wspólne rozdrobnienie
spełnia żądaną nierówność.
Odwrotnie, dla każdego $P$ mamy
$0\le\overline I(f)-\underline I(f)
\le U(f,P)-L(f,P)$; dowolność $\varepsilon$
kończy dowód.
\end{proof}

Dla nieujemnej funkcji dwóch zmiennych sumy
Darboux mają interpretację objętości sum
prostopadłościanów pod wykresem i nad nim.
Dla funkcji ze znakiem są to objętości ze znakiem,
nie zaś objętości geometryczne.

\section{Sumy Riemanna i granica przy zagęszczaniu podziałów}

\begin{definicja}
\label{ma2:def:riemann-box-sum}
Dla każdej komórki wybieramy znacznik
$\xi_\alpha\in Q_\alpha$. Sumą Riemanna nazywamy
\[
 S(f,P,\xi)=\sum_\alpha f(\xi_\alpha)|Q_\alpha|.
\]
Istnienie granicy $I$ przy $\delta(P)\to0$ oznacza,
że dla każdego $\varepsilon>0$ istnieje $\eta>0$
takie, że dla wszystkich podziałów $P$ o
$\delta(P)<\eta$ i wszystkich wyborów znaczników
mamy $|S(f,P,\xi)-I|<\varepsilon$.
\end{definicja}

\begin{twierdzenie}
\label{ma2:thm:riemann-darboux-box}
Dla ograniczonej funkcji $f:Q\to\R$
całkowalność w sensie definicji
\ref{ma2:def:darboux-box} jest równoważna
istnieniu skończonej granicy sum Riemanna
z definicji \ref{ma2:def:riemann-box-sum}.
Obie wartości są równe.
\end{twierdzenie}
\begin{proof}
Niech $|f|\le M$. Gdy $M=0$, wszystkie sumy
są zerowe, więc możemy przyjąć $M>0$.
Załóżmy całkowalność i ustalmy $\varepsilon>0$.
Wybierzmy podział $T$ spełniający
$U(f,T)-L(f,T)<\varepsilon/2$.
Niech $P$ będzie dowolnym podziałem. Nazwijmy
komórkę $P$ złą, jeśli przecina hiperpłaszczyznę
$x_j=t$, gdzie $t$ jest wewnętrznym węzłem $T$
w $j$-tej współrzędnej. Niech $r_j$ oznacza
liczbę tych węzłów. W danej współrzędnej
każda komórka przecinająca $x_j=t$ ma swój
$j$-ty przedział zawarty w
$[t-\delta(P),t+\delta(P)]$.
Suma długości wszystkich takich przedziałów
nie przekracza $2\delta(P)$, bo ich wnętrza
są rozłączne. Sumując po pozostałych
współrzędnych i po węzłach, otrzymujemy
\[
 \sum_{\alpha:\,Q_\alpha\ \text{zła}}|Q_\alpha|
 \le C_T\delta(P),\qquad
 C_T=2\sum_{j=1}^n r_j\prod_{i\ne j}(b_i-a_i).
\]
Zastosowano tu oszacowanie sumy po sumie
skończenie wielu rodzin komórek przez sumę
oszacowań tych rodzin; wielokrotne zliczenie
nie szkodzi, bo objętości są nieujemne.

Niech $R$ będzie wspólnym rozdrobnieniem $P,T$.
Dobre komórki $P$ nie są dzielone przez $T$,
więc ich wkłady do sum dolnych i górnych
pozostają niezmienione. Na złej komórce
różnica między starą i nową sumą dolną
oraz między starą i nową sumą górną
nie przekracza $2M|Q_\alpha|$.
Wobec tego
\[
 L(f,P)\ge L(f,R)-2MC_T\delta(P),\qquad
 U(f,P)\le U(f,R)+2MC_T\delta(P).
\]
Ponieważ $R$ rozdrabnia $T$,
$L(f,T)\le L(f,R)\le U(f,R)\le U(f,T)$.
Ponadto $L(f,P)\le S(f,P,\xi)\le U(f,P)$.
Całka $I_Q(f)$ i obie sumy dla $T$ różnią się
o mniej niż $\varepsilon/2$, zatem
\[
 |S(f,P,\xi)-I_Q(f)|
 <\varepsilon/2+2MC_T\delta(P).
\]
Jeżeli $C_T>0$, wystarczy
$\delta(P)<\varepsilon/(4MC_T)$; jeżeli
$C_T=0$, ostatni składnik jest zerowy.
Oszacowanie nie zależy od znaczników.

Odwrotnie, załóżmy istnienie granicy $I$.
Dla dowolnego $\varepsilon>0$ wybieramy
podział $P$ o tak małej średnicy, aby każda
suma Riemanna różniła się od $I$ o mniej
niż $\varepsilon/4$. Z definicji kresów
w każdej komórce możemy wybrać znacznik
$\xi_\alpha^-$ z
$f(\xi_\alpha^-)<m_\alpha+\varepsilon/(4|Q|)$
oraz $\xi_\alpha^+$ z
$f(\xi_\alpha^+)>M_\alpha-\varepsilon/(4|Q|)$.
Stąd
\[
 U(f,P)-L(f,P)
 < S(f,P,\xi^+)-S(f,P,\xi^-)+\varepsilon/2
 <\varepsilon.
\]
Kryterium Darboux daje całkowalność.
Udowodniona pierwsza implikacja daje granicę
$I_Q(f)$; jednoznaczność granicy (dla dowolnego
ciągu podziałów o średnicy dążącej do zera)
prowadzi do $I=I_Q(f)$.
\end{proof}

\section{Ciągłość i elementarne własności całki}

\begin{twierdzenie}
\label{ma2:thm:riemann-box-continuous}
Każda funkcja ciągła $f:Q\to\R$ jest
całkowalna w sensie Riemanna.
\end{twierdzenie}
\begin{proof}
Zwartość $Q$ daje ograniczoność i jednostajną
ciągłość $f$. Dla $\varepsilon>0$ wybieramy
$\rho>0$ tak, by
$|f(x)-f(y)|<\varepsilon/(2|Q|)$ dla
$x,y\in Q$ o $\|x-y\|<\rho$.
Podział z $\sqrt n\,\delta(P)<\rho$ spełnia
$M_\alpha-m_\alpha\le\varepsilon/(2|Q|)$
dla każdej komórki, nawet gdy kresy nie są osiągane
(wystarczy przejść do kresów w oszacowaniu
różnic). Po przemnożeniu przez objętości
i zsumowaniu dostajemy
$U(f,P)-L(f,P)\le\varepsilon/2$.
\end{proof}

\begin{twierdzenie}
\label{ma2:thm:riemann-box-algebra}
Funkcje całkowalne w sensie Riemanna na $Q$
tworzą algebrę nad $\R$, zawierającą funkcje
stałe i domkniętą na operację $f\mapsto |f|$.
Całka $I_Q$ jest liniowa i monotoniczna, a
\[
 |I_Q(f)|\le I_Q(|f|)\le |Q|\,\|f\|_\infty.
\]
\end{twierdzenie}
\begin{proof}
Dla komórki $C$ oznaczmy oscylację funkcji
$h$ przez $\operatorname{osc}_C h=\sup_C h-\inf_C h$.
Jest ona równa $\sup_{x,y\in C}|h(x)-h(y)|$:
jedna nierówność jest natychmiastowa z definicji
kresów, a druga wynika z wyboru wartości dowolnie
bliskich obu kresom.
Dla ograniczonych $f,g$ z nierówności trójkąta
oraz rozkładu
$f(x)g(x)-f(y)g(y)=f(x)(g(x)-g(y))
+g(y)(f(x)-f(y))$ otrzymujemy
\[
 \begin{aligned}
 \operatorname{osc}_C(\alpha f+\beta g)
 &\le|\alpha|\operatorname{osc}_Cf+
      |\beta|\operatorname{osc}_Cg,\\
 \operatorname{osc}_C(fg)
 &\le\|f\|_\infty\operatorname{osc}_Cg+
      \|g\|_\infty\operatorname{osc}_Cf,\\
 \operatorname{osc}_C|f|&\le\operatorname{osc}_Cf.
 \end{aligned}
\]
Po wybraniu dla $f,g$ podziałów o dowolnie
małych różnicach sum i przejściu do wspólnego
rozdrobnienia kryterium Darboux daje
całkowalność wszystkich wymienionych funkcji.
Funkcje stałe mają równe sumy dolne i górne.
Liniowość wynika z dokładnej liniowości
sum Riemanna dla wspólnych znaczników
i twierdzenia \ref{ma2:thm:riemann-darboux-box}.
Jeśli $f\le g$, nierówność zachodzi dla
wszystkich sum o tych samych znacznikach,
więc zachodzi także dla granic.
Stosując monotoniczność do
$-|f|\le f\le |f|\le\|f\|_\infty$ i używając
liniowości, uzyskujemy żądane oszacowanie.
\end{proof}

\begin{przyklad}[Iloczyn dwóch funkcji współrzędnościowych]
\label{ma2:ex:riemann-box-xy}
Na $Q=[0,1]^2$ funkcja $f(x,y)=xy$ jest ciągła.
Dzielimy każdą oś na $N$ równych odcinków.
Na komórce o indeksach $i,j$ funkcja jest
niemalejąca w obu współrzędnych, więc
\[
 L(f,P_N)=\frac{1}{N^4}
 \sum_{i,j=1}^N(i-1)(j-1)=\frac{(N-1)^2}{4N^2},
 \qquad
 U(f,P_N)=\frac{(N+1)^2}{4N^2}.
\]
Druga tożsamość wynika analogicznie z sumy
$\sum_{i,j}ij=(N(N+1)/2)^2$.
Różnica wynosi $1/N$, a obie sumy dążą
do $1/4$. Zatem $I_Q(f)=1/4$.
Nie użyliśmy tu twierdzenia Fubiniego:
rozkład sumy skończonej jest czysto algebraiczny.
\end{przyklad}

% Drugi blok: kryterium Lebesgue'a i zgodność całek
% w wymiarze n, addytywność, zawartość Jordana i przykład
% geometryczny. Dopiero potem audyty stabilizacyjne.
